\documentclass[11pt,a4paper]{article}

\usepackage[utf8]{inputenc}
\usepackage[T1]{fontenc}
\usepackage[spanish,english]{babel}
\usepackage{amsmath, amssymb, amsfonts}
\usepackage{geometry}
\usepackage{booktabs}
\usepackage{hyperref}
\usepackage{graphicx}
\usepackage{float}
\usepackage[ruled,vlined,linesnumbered]{algorithm2e}

\geometry{margin=2.5cm}

\title{Análisis PRAM y Métricas de un KNN Paralelo con MPI}
\author{
Ian González \and Eduardo Aragon \and Jorge Quenta
}
\date{Mayo 2026}

\begin{document}

\maketitle

\begin{abstract}
Este documento presenta el diseño y análisis teórico de una versión paralela del algoritmo
K-Nearest Neighbors (KNN) usando un modelo de memoria distribuida con MPI. El algoritmo
particiona la data de entrenamiento entre los procesos, replica la data de test mediante una
difusión en árbol, calcula localmente los mejores $K$ vecinos por proceso y combina los resultados
mediante una reducción binaria de listas Top-$K$. Se derivan las expresiones teóricas de tiempo
de cómputo, tiempo de comunicación, speedup, eficiencia, FLOPs y escalabilidad fuerte y débil.
\end{abstract}

\section{Introducción}

K-Nearest Neighbors (KNN) es un algoritmo supervisado de clasificación que asigna a un punto
de prueba la etiqueta mayoritaria entre sus $K$ vecinos más cercanos en el conjunto de
entrenamiento. Dado un conjunto de entrenamiento con $n_{tr}$ puntos, un conjunto de prueba
con $n_{te}$ puntos y una dimensión $d$, el costo principal del algoritmo secuencial se encuentra
en el cálculo de distancias entre todos los puntos de prueba y todos los puntos de entrenamiento.

La versión paralela propuesta distribuye el conjunto de entrenamiento entre $p$ procesos. Cada
proceso recibe aproximadamente $n_{tr}/p$ puntos de entrenamiento, mientras que todos los
procesos reciben el conjunto completo de prueba. Luego, cada proceso calcula los mejores
$K$ vecinos locales para cada punto de prueba. Finalmente, los resultados locales se combinan
mediante una reducción binaria, de manera que el proceso raíz obtiene los mejores $K$ vecinos
globales para cada punto de prueba y realiza la votación mayoritaria.

\section{Notación}

\begin{table}[H]
\centering
\begin{tabular}{ll}
\toprule
Símbolo & Descripción \\
\midrule
$n_{tr}$ & Número de puntos de entrenamiento \\
$n_{te}$ & Número de puntos de prueba \\
$d$ & Número de dimensiones/features \\
$p$ & Número de procesos MPI \\
$s$ & Identificador del proceso, con $1 \leq s \leq p$ \\
$K$ & Número de vecinos considerados por KNN \\
$\ell$ & Número de puntos de entrenamiento por proceso, $\ell = n_{tr}/p$ \\
$\alpha$ & Latencia de envío de un mensaje \\
$\beta$ & Tiempo de transmisión por byte \\
$B_x$ & Bytes por coordenada de un punto \\
$B_c$ & Bytes por candidato Top-$K$ \\
\bottomrule
\end{tabular}
\caption{Notación utilizada en el análisis.}
\end{table}

En el análisis se asume que los datos de entrada son de tipo \texttt{float}, por lo que
$B_x = 4$ bytes. Además, se representa cada candidato como un par
\[
(\text{distancia}, \text{índice})
\]
donde la distancia es un \texttt{float} y el índice es un \texttt{int}. Por tanto,
\[
B_c = 8 \text{ bytes}.
\]
Si se utiliza \texttt{double} o una estructura con padding, $B_c$ puede ser mayor.

\section{Diseño paralelo}

\subsection{Idea general}

El algoritmo usa una descomposición por datos de entrenamiento. Cada proceso recibe un bloque
del conjunto de entrenamiento y calcula, para cada punto de prueba, los mejores $K$ vecinos dentro
de su bloque local. Luego, las listas locales Top-$K$ se reducen en forma de árbol binario.

La operación de combinación entre dos listas Top-$K$ se define como:
\[
\text{MergeTopK}(A,B,K) = \text{mejores } K \text{ candidatos de } A \cup B.
\]

Esta operación permite una reducción porque los mejores $K$ vecinos globales deben estar dentro
de la unión de los mejores $K$ vecinos locales de cada proceso.

\section{PRAM del algoritmo}

\begin{algorithm}[H]
\caption{KNN paralelo por proceso con reducción binaria de Top-$K$}
\KwIn{
$X_{tr}[1,\dots,n_{tr}][1,\dots,d]$, $Y_{tr}[1,\dots,n_{tr}]$,
$X_{te}[1,\dots,n_{te}][1,\dots,d]$,
$p = 2^q$ procesos, proceso $s$ con $1 \leq s \leq p$, número de vecinos $K$
}
\KwOut{
Predicciones $\mathrm{Pred}[1,\dots,n_{te}]$ almacenadas en el proceso $s=1$
}

\BlankLine
$\ell \leftarrow n_{tr}/p$\;

\BlankLine
\tcp{1. Distribución del conjunto de entrenamiento}
\If{$s = 1$}{
    \For{$r \leftarrow 1$ \KwTo $p$}{
        enviar $X_{tr}[\ell(r-1)+1,\dots,\ell r][1,\dots,d]$ al proceso $r$\;
        enviar $Y_{tr}[\ell(r-1)+1,\dots,\ell r]$ al proceso $r$\;
    }
}
recibir $\mathrm{LocalX}[1,\dots,\ell][1,\dots,d]$\;
recibir $\mathrm{LocalY}[1,\dots,\ell]$\;

\BlankLine
\tcp{2. Difusión en árbol del conjunto de prueba}
\For{$h \leftarrow 1$ \KwTo $\log_2(p)$}{
    $\mathrm{half} \leftarrow 2^{h-1}$\;
    \If{$s \leq \mathrm{half}$}{
        $\mathrm{dest} \leftarrow s + \mathrm{half}$\;
        \If{$\mathrm{dest} \leq p$}{
            enviar $X_{te}[1,\dots,n_{te}][1,\dots,d]$ a $\mathrm{dest}$\;
        }
    }
    \ElseIf{$\mathrm{half} < s \leq 2^h$}{
        $\mathrm{src} \leftarrow s - \mathrm{half}$\;
        recibir $X_{te}[1,\dots,n_{te}][1,\dots,d]$ desde $\mathrm{src}$\;
    }
}

\BlankLine
\tcp{3. Cálculo local de Top-$K$}
\For{$t \leftarrow 1$ \KwTo $n_{te}$}{
    $\mathrm{TopK}_s[t] \leftarrow \emptyset$\;
    
    \For{$i \leftarrow 1$ \KwTo $\ell$}{
        $\mathrm{dist} \leftarrow \mathrm{distance}(X_{te}[t], \mathrm{LocalX}[i])$\;
        
        $\mathrm{idx} \leftarrow \ell(s-1)+i$\;
        
        $\mathrm{candidate} \leftarrow (\mathrm{dist}, \mathrm{idx})$\;
        
        $\mathrm{TopK}_s[t] \leftarrow 
        \mathrm{InsertAndKeepBestK}(\mathrm{TopK}_s[t], \mathrm{candidate}, K)$\;
    }
}

\BlankLine
\tcp{4. Reducción binaria de Top-$K$}
\For{$h \leftarrow 1$ \KwTo $\log_2(p)$}{
    $\mathrm{block\_size} \leftarrow 2^h$\;
    $\mathrm{half} \leftarrow 2^{h-1}$\;
    
    \BlankLine
    \If{$(s-1) \bmod \mathrm{block\_size} = \mathrm{half}$}{
        $\mathrm{parent} \leftarrow s - \mathrm{half}$\;
        
        enviar $\mathrm{TopK}_s[1,\dots,n_{te}]$ a $\mathrm{parent}$\;
        
        \textbf{break}\;
    }
    
    \BlankLine
    \ElseIf{$(s-1) \bmod \mathrm{block\_size} = 0$}{
        $\mathrm{child} \leftarrow s + \mathrm{half}$\;
        
        \If{$\mathrm{child} \leq p$}{
            recibir $\mathrm{TopK}_{child}[1,\dots,n_{te}]$ desde $\mathrm{child}$\;
            
            \For{$t \leftarrow 1$ \KwTo $n_{te}$}{
                $\mathrm{TopK}_s[t] \leftarrow 
                \mathrm{MergeTopK}(\mathrm{TopK}_s[t], \mathrm{TopK}_{child}[t], K)$\;
            }
        }
    }
}
\end{algorithm}

\begin{algorithm}
\tcp{5. Votación mayoritaria en el proceso raíz}
\If{$s = 1$}{
    \For{$t \leftarrow 1$ \KwTo $n_{te}$}{
        $\mathrm{Pred}[t] \leftarrow \mathrm{MajorityVote}(\mathrm{TopK}_s[t], Y_{tr})$\;
    }
}
\end{algorithm}

\subsection{Interpretación de la reducción}

En la fase de reducción, cada proceso activo representa el acumulador del subárbol izquierdo.
Por ello, en cada nivel $h$, solo se envía la información del subárbol derecho hacia el proceso
padre. El proceso padre fusiona su $\mathrm{TopK}_s$ actual con el $\mathrm{TopK}$ recibido y conserva
los mejores $K$ candidatos para cada punto de prueba.

La condición de envío es:
\[
(s-1) \bmod 2^h = 2^{h-1},
\]
mientras que la condición de recepción es:
\[
(s-1) \bmod 2^h = 0.
\]

Para $p=8$, los envíos son:
\[
h=1: \quad 2\to1,\;4\to3,\;6\to5,\;8\to7,
\]
\[
h=2: \quad 3\to1,\;7\to5,
\]
\[
h=3: \quad 5\to1.
\]

Al finalizar, el proceso $s=1$ contiene el Top-$K$ global para cada punto de prueba.

\section{Complejidad de cómputo}

\subsection{Cálculo local de distancias}

Cada proceso compara sus $\ell = n_{tr}/p$ puntos de entrenamiento contra los $n_{te}$ puntos
de prueba. Cada distancia euclidiana en dimensión $d$ cuesta:
\[
O(d).
\]

Se asume que la lista Top-$K$ de cada punto de prueba se mantiene con una cola de prioridad
de tamaño máximo $K$. Por tanto, insertar o actualizar un candidato cuesta:
\[
O(\log K).
\]

Así, el costo local principal por proceso es:
\[
T_{\text{local}}
=
O\left(
\frac{n_{tr}}{p} n_{te}(d+\log K)
\right).
\]

\subsection{Cómputo durante la reducción}

En cada nivel de la reducción, un proceso receptor fusiona dos listas Top-$K$ para cada punto
de prueba. Fusionar dos listas de tamaño $K$ cuesta:
\[
O(K).
\]

Como hay $n_{te}$ puntos de prueba, el costo por nivel es:
\[
O(n_{te}K).
\]

La reducción tiene profundidad $\log_2 p$, por lo que:
\[
T_{\text{reduce,comp}}
=
O(n_{te}K\log p).
\]

\subsection{Votación mayoritaria}

En el proceso raíz, la votación mayoritaria se realiza para cada punto de prueba usando sus
$K$ vecinos globales:
\[
T_{\text{vote}}
=
O(n_{te}K).
\]

\subsection{Tiempo total de cómputo}

Por tanto, el tiempo de cómputo paralelo es:
\[
T_{\text{comp}}
=
O\left(
\frac{n_{tr}n_{te}}{p}(d+\log K)
+
n_{te}K\log p
+
n_{te}K
\right).
\]

Como el término de votación suele estar dominado por el cálculo de distancias, puede
simplificarse como:
\[
\boxed{
T_{\text{comp}}
=
O\left(
\frac{n_{tr}n_{te}}{p}(d+\log K)
+
n_{te}K\log p
\right)
}.
\]

\section{Complejidad de comunicación}

El costo de comunicación se modela como:
\[
T_{\text{msg}} = \alpha + x\beta,
\]
donde $\alpha$ es la latencia, $\beta$ es el costo por byte transmitido y $x$ es el número
de bytes enviados.

\subsection{Distribución del conjunto de entrenamiento}

El proceso raíz distribuye aproximadamente $\ell = n_{tr}/p$ puntos a cada proceso. Cada punto
tiene $d$ coordenadas y cada coordenada ocupa $B_x = 4$ bytes. La cantidad total de bytes
enviados es:
\[
4n_{tr}d.
\]

Además, hay $p$ mensajes. Por tanto:
\[
T_{\text{scatter}}
=
O(p\alpha + 4\beta n_{tr}d).
\]

Si se incluyen las etiquetas $Y_{tr}$, debe añadirse:
\[
O(4\beta n_{tr}),
\]
asumiendo etiquetas enteras de 4 bytes. Como usualmente $d \gg 1$, este término es menor.

\subsection{Difusión en árbol del conjunto de prueba}

El conjunto de prueba contiene $n_{te}d$ valores. Como cada valor ocupa 4 bytes:
\[
x = 4n_{te}d.
\]

Usando difusión en árbol, la profundidad es $\log p$, por lo que:
\[
T_{\text{bcast}}
=
O\left(
(\alpha + 4\beta n_{te}d)\log p
\right).
\]

\subsection{Reducción binaria de listas Top-$K$}

Cada proceso envía, para cada punto de prueba, una lista de $K$ candidatos. Cada candidato
contiene distancia e índice, por lo que ocupa $B_c = 8$ bytes. El tamaño del mensaje de
reducción es:
\[
x = 8n_{te}K.
\]

Como la reducción tiene profundidad $\log p$:
\[
T_{\text{reduce,comm}}
=
O\left(
(\alpha + 8\beta n_{te}K)\log p
\right).
\]

\subsection{Tiempo total de comunicación}

Sumando los términos:
\[
\boxed{
T_{\text{comm}}
=
O\left(
p\alpha
+
4\beta n_{tr}d
+
(\alpha + 4\beta n_{te}d)\log p
+
(\alpha + 8\beta n_{te}K)\log p
\right)
}.
\]

\section{Tiempo paralelo total}

El tiempo paralelo total es:
\[
T_p = T_{\text{comp}} + T_{\text{comm}}.
\]

Sustituyendo:
\[
\boxed{
T_p
=
O\left(
\frac{n_{tr}n_{te}}{p}(d+\log K)
+
n_{te}K\log p
+
p\alpha
+
4\beta n_{tr}d
+
(\alpha + 4\beta n_{te}d)\log p
+
(\alpha + 8\beta n_{te}K)\log p
\right)
}.
\]

\section{Tiempo secuencial}

En la versión secuencial, se comparan todos los puntos de prueba contra todos los puntos de
entrenamiento. Con una cola de prioridad de tamaño $K$, se obtiene:
\[
\boxed{
T_1
=
O\left(
n_{tr}n_{te}(d+\log K)
+
n_{te}K
\right)
}.
\]

Como $n_{tr}d$ suele ser mucho mayor que $K$, normalmente se simplifica como:
\[
T_1
=
O(n_{tr}n_{te}(d+\log K)).
\]

\section{Speedup y eficiencia}

\subsection{Speedup teórico}

El speedup se define como:
\[
S_p = \frac{T_1}{T_p}.
\]

Entonces:
\[
\boxed{
S_p
=
\frac{
O(n_{tr}n_{te}(d+\log K))
}{
O\left(
\frac{n_{tr}n_{te}}{p}(d+\log K)
+
n_{te}K\log p
+
T_{\text{comm}}
\right)
}
}.
\]

Si se ignora comunicación y reducción, el speedup ideal sería:
\[
S_p^{ideal} = p.
\]

Sin embargo, en la práctica:
\[
S_p < p,
\]
debido a comunicación, reducción y desbalance de carga.

\subsection{Eficiencia}

La eficiencia se define como:
\[
E_p = \frac{S_p}{p}.
\]

Sustituyendo:
\[
\boxed{
E_p
=
\frac{T_1}{pT_p}
}.
\]

La eficiencia mide qué fracción del rendimiento ideal se aprovecha al usar $p$ procesos.
Cuando $E_p \approx 1$, el paralelismo es casi ideal. Cuando $E_p$ disminuye, significa que
la comunicación, sincronización o el desbalance de carga están dominando.

\section{FLOPs}

\subsection{FLOPs por distancia euclidiana}

Para dos puntos $x,y \in \mathbb{R}^d$, se calcula la distancia euclidiana al cuadrado:
\[
\|x-y\|^2 = \sum_{j=1}^{d}(x_j-y_j)^2.
\]

Para cada dimensión se realizan aproximadamente:
\begin{itemize}
    \item 1 resta: $x_j-y_j$,
    \item 1 multiplicación: $(x_j-y_j)^2$,
    \item 1 suma acumulada.
\end{itemize}

Por tanto, se cuentan aproximadamente:
\[
3d
\]
FLOPs por distancia.

No se calcula la raíz cuadrada, porque para KNN basta comparar distancias al cuadrado; la raíz
cuadrada preserva el orden de las distancias y añade costo innecesario.

\subsection{FLOPs totales de la región paralelizable}

La región paralelizable principal calcula:
\[
n_{tr}n_{te}
\]
distancias. Por tanto:
\[
\boxed{
\mathrm{FLOPs}_{total}
=
3n_{tr}n_{te}d
}.
\]

Por proceso:
\[
\boxed{
\mathrm{FLOPs}_{proc}
=
\frac{3n_{tr}n_{te}d}{p}
}.
\]

\subsection{FLOPs por segundo}

Si $T_{\text{comp}}^{exp}$ es el tiempo experimental de cómputo medido, entonces:
\[
\boxed{
\mathrm{FLOPs/s}
=
\frac{3n_{tr}n_{te}d}{T_{\text{comp}}^{exp}}
}.
\]

Si se desea calcular GFLOPs/s:
\[
\boxed{
\mathrm{GFLOPs/s}
=
\frac{3n_{tr}n_{te}d}{10^9T_{\text{comp}}^{exp}}
}.
\]

\section{Normalización de la complejidad teórica}

La notación Big-$O$ describe crecimiento asintótico, pero no da tiempos en segundos. Para comparar
la teoría con los experimentos, se normaliza la expresión teórica usando una constante de
proporcionalidad.

Se propone ajustar:
\[
T_{\text{theory}}(p)
=
C_1
\frac{n_{tr}n_{te}}{p}(d+\log K)
+
C_2 n_{te}K\log p
+
C_3 T_{\text{comm}}(p),
\]
donde $C_1$, $C_2$ y $C_3$ se estiman a partir de los resultados experimentales.

Una forma simple de normalizar es usar el tiempo experimental para $p=1$:
\[
C_1
=
\frac{T_{exp}(1)}{n_{tr}n_{te}(d+\log K)}.
\]

Luego se grafica:
\[
T_{\text{theory}}(p)
\quad \text{vs.} \quad
T_{\text{exp}}(p).
\]

\section{Escalabilidad fuerte}

La escalabilidad fuerte evalúa cómo cambia el tiempo al aumentar $p$ manteniendo fijo el tamaño
del problema:
\[
n_{tr}, n_{te}, d, K = \text{constantes}.
\]

En este caso:
\[
T_{\text{comp}} =
O\left(
\frac{n_{tr}n_{te}}{p}(d+\log K)
\right)
\]
disminuye al aumentar $p$.

Sin embargo, los términos de comunicación:
\[
p\alpha,
\]
\[
(\alpha + 4\beta n_{te}d)\log p,
\]
\[
(\alpha + 8\beta n_{te}K)\log p
\]
no disminuyen con $p$. Por tanto, para valores pequeños de $p$, el algoritmo debería mostrar
mejoras cercanas a lineales. Para valores grandes de $p$, el costo de comunicación y reducción
empieza a dominar.

El comportamiento esperado es:
\[
S_p \approx p
\quad \text{para } p \text{ pequeño},
\]
pero:
\[
E_p = \frac{S_p}{p}
\]
decrece cuando $p$ crece demasiado.

\subsection{Estimación del número óptimo de procesos}

Una aproximación útil consiste en escribir:
\[
T_p \approx \frac{A}{p} + C\log p + Dp,
\]
donde:
\[
A = n_{tr}n_{te}(d+\log K),
\]
\[
C \approx n_{te}K + \alpha + 4\beta n_{te}d + 8\beta n_{te}K,
\]
y $D$ representa el costo asociado a términos lineales como $p\alpha$.

El número óptimo de procesos se estima experimentalmente como:
\[
\boxed{
p^* = \arg\min_p T_p^{exp}
}.
\]

También puede definirse como el mayor $p$ antes de que la eficiencia caiga bajo un umbral,
por ejemplo:
\[
E_p < 0.5.
\]

\section{Escalabilidad débil}

La escalabilidad débil evalúa qué ocurre cuando aumentan simultáneamente el tamaño del problema
y el número de procesos.

Para este diseño, como se particiona con respecto a $n_{tr}$, una forma natural de escalabilidad
débil es mantener constante la cantidad de entrenamiento por proceso:
\[
\frac{n_{tr}}{p} = \ell = \text{constante}.
\]

Si además se mantiene fijo $n_{te}$, $d$ y $K$, entonces:
\[
T_{\text{comp}}
=
O(\ell n_{te}(d+\log K)),
\]
que es constante con respecto a $p$.

Pero la comunicación aún crece:
\[
T_{\text{bcast}} =
O((\alpha + 4\beta n_{te}d)\log p),
\]
\[
T_{\text{reduce}} =
O((\alpha + 8\beta n_{te}K)\log p).
\]

Por tanto, bajo esta forma de weak scaling:
\[
T_p
=
O(1 + \log p).
\]

Esto implica que el algoritmo puede escalar razonablemente, pero no de forma perfecta, debido
a la comunicación global de $X_{te}$ y la reducción de Top-$K$.

\subsection{Caso donde crecen entrenamiento y prueba}

Si se escala todo el dataset manteniendo fija la proporción entre entrenamiento y prueba:
\[
n_{tr} \propto p,
\quad
n_{te} \propto p,
\]
entonces el trabajo total de KNN crece como:
\[
O(n_{tr}n_{te}d) = O(p^2d).
\]

Al dividir solo $n_{tr}$ entre $p$ procesos:
\[
T_{\text{comp}}
=
O\left(
\frac{n_{tr}n_{te}d}{p}
\right)
=
O(pd).
\]

Por tanto, si tanto $n_{tr}$ como $n_{te}$ crecen proporcionalmente a $p$, el tiempo por proceso
no se mantiene constante, sino que crece linealmente con $p$. Esto no representa weak scaling
ideal, pero sí representa el comportamiento esperado cuando se aumenta el tamaño completo del
dataset.

\section{Métricas experimentales a reportar}

Para cumplir con el análisis experimental, se recomienda medir para cada combinación de
$n$ y $p$:

\begin{itemize}
    \item Tiempo total de ejecución:
    \[
    T_{\text{total}}
    \]
    
    \item Tiempo de cómputo:
    \[
    T_{\text{comp}}^{exp}
    \]
    
    \item Tiempo de comunicación:
    \[
    T_{\text{comm}}^{exp}
    \]
    
    \item Accuracy:
    \[
    \mathrm{Accuracy} =
    \frac{\text{predicciones correctas}}{\text{total de predicciones}}
    \]
    
    \item Speedup:
    \[
    S_p =
    \frac{T_1}{T_p}
    \]
    
    \item Eficiencia:
    \[
    E_p =
    \frac{S_p}{p}
    \]
    
    \item GFLOPs/s:
    \[
    \mathrm{GFLOPs/s}
    =
    \frac{3n_{tr}n_{te}d}{10^9T_{\text{comp}}^{exp}}
    \]
\end{itemize}

\section{Relación con MPI}

El PRAM propuesto puede implementarse en Python con \texttt{mpi4py}. Las operaciones principales
son:

\begin{table}[H]
\centering
\begin{tabular}{ll}
\toprule
Fase del PRAM & Directiva MPI sugerida \\
\midrule
Distribución de entrenamiento & \texttt{comm.scatter} o \texttt{comm.scatterv} \\
Difusión de test & \texttt{comm.bcast} \\
Cálculo local & Código local NumPy/Python \\
Reducción Top-$K$ & \texttt{comm.send}/\texttt{comm.recv} en árbol \\
Alternativa simple & \texttt{comm.gather} y merge en root \\
Votación final & Proceso root \\
\bottomrule
\end{tabular}
\caption{Correspondencia entre el PRAM y una implementación con MPI.}
\end{table}

Aunque la reducción binaria manual es más eficiente que recolectar todos los Top-$K$ en el
proceso raíz, una primera versión beta puede usar \texttt{comm.gather} para simplificar la
implementación. Luego, una versión posterior puede reemplazar ese \texttt{gather} por la reducción
binaria descrita en el PRAM.

\section{Versiones beta sugeridas}

\begin{itemize}
    \item \textbf{Versión beta 1:} implementación secuencial validada con \texttt{load\_digits}.
    Se mide accuracy y tiempo base $T_1$.
    
    \item \textbf{Versión beta 2:} implementación MPI con \texttt{comm.bcast}, \texttt{comm.scatter}
    y \texttt{comm.gather}. Cada proceso calcula Top-$K$ local y el root combina todos los
    resultados.
    
    \item \textbf{Versión beta 3:} reemplazo de \texttt{comm.gather} por reducción binaria de
    Top-$K$, reduciendo el cuello de botella en el proceso raíz.
\end{itemize}

% =====================================================================
%  Fragmento integrable en main.tex mediante:  % =====================================================================
%  Fragmento integrable en main.tex mediante:  % =====================================================================
%  Fragmento integrable en main.tex mediante:  \input{generated/isoeficiencia.tex}
%  Reutiliza la notación, los paquetes (amsmath, booktabs, float) y las
%  complejidades ya derivadas en el documento principal. No define preámbulo.
% =====================================================================

\subsection{Análisis de isoeficiencia (condición de escalabilidad)}

En este curso la \textbf{isoeficiencia} se evalúa como una \emph{condición de escalabilidad
débil} (weak scaling) basada en la Ley de Gustafson. No se busca una función de isoeficiencia
aislada, sino la relación que dicta cómo debe crecer el tamaño del problema $n$ en función del
número de procesos $p$ para que la eficiencia $E$ se mantenga constante ($E=O(1)$). Formalizamos
aquí la condición que ya se anticipó de manera cualitativa en las secciones de escalabilidad
fuerte y débil, y comparamos el paradigma usado (MPI) con el de memoria compartida (OMP/PRAM).

\subsubsection{Función de overhead y forma canónica de la eficiencia}

Para obtener una única condición de escalabilidad tomamos un tamaño de problema $n$ con
$n_{tr}=n_{te}=n$ (el dataset completo escala). De las complejidades ya derivadas, el trabajo
total (tiempo secuencial dominante, omitiendo $\log K$) es
\[
W \equiv T_1 = O(n^2 d),
\]
y el cómputo paralelo reparte ese trabajo idealmente, $T_{\text{comp}} = O(n^2 d / p)$.

Definimos la \textbf{función de overhead total} $T_o$ como el trabajo que no realiza el algoritmo
secuencial:
\[
T_o(p) = p\,T_p - W = p\,T_{\text{comp}} + p\,T_{\text{comm}} - W.
\]
Como $p\,T_{\text{comp}} = W$, todo el overhead proviene de la comunicación y la sincronización:
\[
\boxed{\; T_o(p) = p\, T_{\text{comm}}. \;}
\]
Sustituyendo en la eficiencia se obtiene la forma canónica del curso:
\[
E = \frac{S_p}{p} = \frac{T_1}{p\,T_p} = \frac{W}{W + T_o}
  = \frac{1}{\,1 + \dfrac{T_o}{W}\,}.
\]
Para que $E=O(1)$ sea constante, el término de overhead $T_o/W$ debe permanecer constante: esa
igualdad, despejada, define la condición de isoeficiencia.

\subsubsection{Isoeficiencia en MPI (memoria distribuida)}

En MPI la comunicación es explícita y es el cuello de botella. Sustituyendo $n_{tr}=n_{te}=n$ en
$T_{\text{comm}}$ y conservando los términos dominantes:
\[
T_o = p\,T_{\text{comm}}
    \approx \underbrace{p^2\alpha}_{(1)\ \text{latencia}}
    \;+\; \underbrace{p\,\beta\, n\, d}_{(2)\ \text{scatter}}
    \;+\; \underbrace{p\,\beta\, n\, d\,\log p}_{(3)\ \text{bcast/reducción}}.
\]
El término (3) proviene de difundir el conjunto de prueba ($O(nd)$ bytes) y de reducir las listas
Top-$K$ a lo largo de una jerarquía de profundidad $\log p$; domina sobre (2). Imponemos
$T_o/W=\text{const}$ con $W=n^2 d$:
\[
(1)\quad \frac{p^2\alpha}{n^2 d}=\text{const} \;\Rightarrow\; n \propto p,
\qquad
(2)\quad \frac{p\,\beta n d}{n^2 d}=\frac{p\beta}{n}=\text{const} \;\Rightarrow\; n \propto p,
\]
\[
(3)\quad \frac{p\,\beta n d\,\log p}{n^2 d}=\frac{p\beta\log p}{n}=\text{const}
\;\Rightarrow\; n \propto p\,\log p.
\]
La condición más restrictiva es la del término (3). Por tanto:
\[
\boxed{\;
\text{Isoeficiencia MPI:}\quad n \propto p\,\log p,
\qquad
W_{iso} = \Theta\!\big(p^2 \log^2 p\big).
\;}
\]
Es decir, el dataset debe crecer ligeramente más rápido que $p$ (factor $\log p$) para amortizar
la difusión del test y la reducción global de candidatos sobre la red. Si la reducción se hiciera
con un \texttt{gather} plano al proceso raíz en vez de la reducción en árbol, el término dominante
sería $O(p)$ en lugar de $O(\log p)$, saturando el enlace del root y empeorando notablemente la
isoeficiencia.

\subsubsection{Isoeficiencia en OMP / PRAM (memoria compartida)}

En OpenMP los hilos comparten el mismo espacio de direcciones (modelo UMA), lo que cambia el
overhead de forma fundamental:
\begin{itemize}
    \item \textbf{No hay difusión de datos:} el conjunto de prueba es visible por todos los hilos,
    por lo que \emph{desaparecen} los términos $p\alpha$ y $\beta n d$ de la comunicación
    explícita ($T_{\text{comm}} \approx 0$).
    \item El único overhead proviene de la \textbf{sincronización de hilos} (barreras) y de la
    reducción de listas Top-$K$ de profundidad $O(\log p)$, realizada mediante accesos directos a
    memoria.
\end{itemize}
Modelando ese overhead (con $K$ constante):
\[
T_o^{\text{omp}} \approx p\cdot O(n_{te}K\log p) + O(p)
= O(p\,n\,\log p) + O(p).
\]
Imponiendo $T_o^{\text{omp}}/W=\text{const}$ con $W=n^2 d$, y observando que ya no aparece el
factor de transporte de datos $\beta d$:
\[
\frac{p\,n\,\log p}{n^2 d}=\frac{p\log p}{n\,d}=\text{const}
\;\Rightarrow\;
n \propto \frac{p\,\log p}{d}.
\]
\[
\boxed{\;
\text{Isoeficiencia OMP:}\quad n \propto \frac{p\,\log p}{d}\ \ (\text{a lo sumo } n\propto p),
\qquad
W_{iso} = \Theta(p^2)\ \text{sin la constante de red } \beta.
\;}
\]
OMP tiene una isoeficiencia asintóticamente mejor (no paga latencia ni ancho de banda de red),
\textbf{pero} $p$ está estrictamente acotado al número de núcleos físicos del nodo y todo el
dataset debe caber en su RAM.

\subsubsection{Comparación y elección del paradigma}

\begin{table}[H]
\centering
\begin{tabular}{lll}
\toprule
Criterio & MPI (memoria distribuida) & OMP (memoria compartida) \\
\midrule
Comunicación $T_{\text{comm}}$ & Explícita: $\alpha,\beta$ (red) & $\approx 0$ (acceso directo) \\
Overhead dominante & Difusión/reducción $\beta n d\log p$ & Sincronización (barreras) \\
Condición de isoeficiencia & $n \propto p\log p$ & $n \propto p\log p / d$ (más suave) \\
Trabajo $W_{iso}$ & $\Theta(p^2\log^2 p)$ & $\Theta(p^2)$ sin $\beta$ \\
Límite de $p$ & Horizontal (varios nodos) & Núcleos de un solo nodo \\
Límite de $n$ & Suma de RAM de los nodos & RAM de un solo nodo \\
\bottomrule
\end{tabular}
\caption{Comparación de la escalabilidad de ambos paradigmas para el KNN.}
\label{tab:iso-comparacion}
\end{table}

Aunque OMP escala mejor asintóticamente, su paralelismo está limitado a los núcleos y la memoria
de un solo nodo. En este proyecto se eligió \textbf{MPI por un tema de memoria}: al distribuir el
conjunto de entrenamiento, el mismo código puede superar el límite de RAM de un nodo y escalar
$n$ y $p$ horizontalmente sin reescritura. Los experimentos evaluaron hasta $p=32$ procesos, y la
elección del paradigma de memoria distribuida es la que habilita crecer más allá de ese punto
cuando el dataset no cabe en un único nodo.


%  Reutiliza la notación, los paquetes (amsmath, booktabs, float) y las
%  complejidades ya derivadas en el documento principal. No define preámbulo.
% =====================================================================

\subsection{Análisis de isoeficiencia (condición de escalabilidad)}

En este curso la \textbf{isoeficiencia} se evalúa como una \emph{condición de escalabilidad
débil} (weak scaling) basada en la Ley de Gustafson. No se busca una función de isoeficiencia
aislada, sino la relación que dicta cómo debe crecer el tamaño del problema $n$ en función del
número de procesos $p$ para que la eficiencia $E$ se mantenga constante ($E=O(1)$). Formalizamos
aquí la condición que ya se anticipó de manera cualitativa en las secciones de escalabilidad
fuerte y débil, y comparamos el paradigma usado (MPI) con el de memoria compartida (OMP/PRAM).

\subsubsection{Función de overhead y forma canónica de la eficiencia}

Para obtener una única condición de escalabilidad tomamos un tamaño de problema $n$ con
$n_{tr}=n_{te}=n$ (el dataset completo escala). De las complejidades ya derivadas, el trabajo
total (tiempo secuencial dominante, omitiendo $\log K$) es
\[
W \equiv T_1 = O(n^2 d),
\]
y el cómputo paralelo reparte ese trabajo idealmente, $T_{\text{comp}} = O(n^2 d / p)$.

Definimos la \textbf{función de overhead total} $T_o$ como el trabajo que no realiza el algoritmo
secuencial:
\[
T_o(p) = p\,T_p - W = p\,T_{\text{comp}} + p\,T_{\text{comm}} - W.
\]
Como $p\,T_{\text{comp}} = W$, todo el overhead proviene de la comunicación y la sincronización:
\[
\boxed{\; T_o(p) = p\, T_{\text{comm}}. \;}
\]
Sustituyendo en la eficiencia se obtiene la forma canónica del curso:
\[
E = \frac{S_p}{p} = \frac{T_1}{p\,T_p} = \frac{W}{W + T_o}
  = \frac{1}{\,1 + \dfrac{T_o}{W}\,}.
\]
Para que $E=O(1)$ sea constante, el término de overhead $T_o/W$ debe permanecer constante: esa
igualdad, despejada, define la condición de isoeficiencia.

\subsubsection{Isoeficiencia en MPI (memoria distribuida)}

En MPI la comunicación es explícita y es el cuello de botella. Sustituyendo $n_{tr}=n_{te}=n$ en
$T_{\text{comm}}$ y conservando los términos dominantes:
\[
T_o = p\,T_{\text{comm}}
    \approx \underbrace{p^2\alpha}_{(1)\ \text{latencia}}
    \;+\; \underbrace{p\,\beta\, n\, d}_{(2)\ \text{scatter}}
    \;+\; \underbrace{p\,\beta\, n\, d\,\log p}_{(3)\ \text{bcast/reducción}}.
\]
El término (3) proviene de difundir el conjunto de prueba ($O(nd)$ bytes) y de reducir las listas
Top-$K$ a lo largo de una jerarquía de profundidad $\log p$; domina sobre (2). Imponemos
$T_o/W=\text{const}$ con $W=n^2 d$:
\[
(1)\quad \frac{p^2\alpha}{n^2 d}=\text{const} \;\Rightarrow\; n \propto p,
\qquad
(2)\quad \frac{p\,\beta n d}{n^2 d}=\frac{p\beta}{n}=\text{const} \;\Rightarrow\; n \propto p,
\]
\[
(3)\quad \frac{p\,\beta n d\,\log p}{n^2 d}=\frac{p\beta\log p}{n}=\text{const}
\;\Rightarrow\; n \propto p\,\log p.
\]
La condición más restrictiva es la del término (3). Por tanto:
\[
\boxed{\;
\text{Isoeficiencia MPI:}\quad n \propto p\,\log p,
\qquad
W_{iso} = \Theta\!\big(p^2 \log^2 p\big).
\;}
\]
Es decir, el dataset debe crecer ligeramente más rápido que $p$ (factor $\log p$) para amortizar
la difusión del test y la reducción global de candidatos sobre la red. Si la reducción se hiciera
con un \texttt{gather} plano al proceso raíz en vez de la reducción en árbol, el término dominante
sería $O(p)$ en lugar de $O(\log p)$, saturando el enlace del root y empeorando notablemente la
isoeficiencia.

\subsubsection{Isoeficiencia en OMP / PRAM (memoria compartida)}

En OpenMP los hilos comparten el mismo espacio de direcciones (modelo UMA), lo que cambia el
overhead de forma fundamental:
\begin{itemize}
    \item \textbf{No hay difusión de datos:} el conjunto de prueba es visible por todos los hilos,
    por lo que \emph{desaparecen} los términos $p\alpha$ y $\beta n d$ de la comunicación
    explícita ($T_{\text{comm}} \approx 0$).
    \item El único overhead proviene de la \textbf{sincronización de hilos} (barreras) y de la
    reducción de listas Top-$K$ de profundidad $O(\log p)$, realizada mediante accesos directos a
    memoria.
\end{itemize}
Modelando ese overhead (con $K$ constante):
\[
T_o^{\text{omp}} \approx p\cdot O(n_{te}K\log p) + O(p)
= O(p\,n\,\log p) + O(p).
\]
Imponiendo $T_o^{\text{omp}}/W=\text{const}$ con $W=n^2 d$, y observando que ya no aparece el
factor de transporte de datos $\beta d$:
\[
\frac{p\,n\,\log p}{n^2 d}=\frac{p\log p}{n\,d}=\text{const}
\;\Rightarrow\;
n \propto \frac{p\,\log p}{d}.
\]
\[
\boxed{\;
\text{Isoeficiencia OMP:}\quad n \propto \frac{p\,\log p}{d}\ \ (\text{a lo sumo } n\propto p),
\qquad
W_{iso} = \Theta(p^2)\ \text{sin la constante de red } \beta.
\;}
\]
OMP tiene una isoeficiencia asintóticamente mejor (no paga latencia ni ancho de banda de red),
\textbf{pero} $p$ está estrictamente acotado al número de núcleos físicos del nodo y todo el
dataset debe caber en su RAM.

\subsubsection{Comparación y elección del paradigma}

\begin{table}[H]
\centering
\begin{tabular}{lll}
\toprule
Criterio & MPI (memoria distribuida) & OMP (memoria compartida) \\
\midrule
Comunicación $T_{\text{comm}}$ & Explícita: $\alpha,\beta$ (red) & $\approx 0$ (acceso directo) \\
Overhead dominante & Difusión/reducción $\beta n d\log p$ & Sincronización (barreras) \\
Condición de isoeficiencia & $n \propto p\log p$ & $n \propto p\log p / d$ (más suave) \\
Trabajo $W_{iso}$ & $\Theta(p^2\log^2 p)$ & $\Theta(p^2)$ sin $\beta$ \\
Límite de $p$ & Horizontal (varios nodos) & Núcleos de un solo nodo \\
Límite de $n$ & Suma de RAM de los nodos & RAM de un solo nodo \\
\bottomrule
\end{tabular}
\caption{Comparación de la escalabilidad de ambos paradigmas para el KNN.}
\label{tab:iso-comparacion}
\end{table}

Aunque OMP escala mejor asintóticamente, su paralelismo está limitado a los núcleos y la memoria
de un solo nodo. En este proyecto se eligió \textbf{MPI por un tema de memoria}: al distribuir el
conjunto de entrenamiento, el mismo código puede superar el límite de RAM de un nodo y escalar
$n$ y $p$ horizontalmente sin reescritura. Los experimentos evaluaron hasta $p=32$ procesos, y la
elección del paradigma de memoria distribuida es la que habilita crecer más allá de ese punto
cuando el dataset no cabe en un único nodo.


%  Reutiliza la notación, los paquetes (amsmath, booktabs, float) y las
%  complejidades ya derivadas en el documento principal. No define preámbulo.
% =====================================================================

\subsection{Análisis de isoeficiencia (condición de escalabilidad)}

En este curso la \textbf{isoeficiencia} se evalúa como una \emph{condición de escalabilidad
débil} (weak scaling) basada en la Ley de Gustafson. No se busca una función de isoeficiencia
aislada, sino la relación que dicta cómo debe crecer el tamaño del problema $n$ en función del
número de procesos $p$ para que la eficiencia $E$ se mantenga constante ($E=O(1)$). Formalizamos
aquí la condición que ya se anticipó de manera cualitativa en las secciones de escalabilidad
fuerte y débil, y comparamos el paradigma usado (MPI) con el de memoria compartida (OMP/PRAM).

\subsubsection{Función de overhead y forma canónica de la eficiencia}

Para obtener una única condición de escalabilidad tomamos un tamaño de problema $n$ con
$n_{tr}=n_{te}=n$ (el dataset completo escala). De las complejidades ya derivadas, el trabajo
total (tiempo secuencial dominante, omitiendo $\log K$) es
\[
W \equiv T_1 = O(n^2 d),
\]
y el cómputo paralelo reparte ese trabajo idealmente, $T_{\text{comp}} = O(n^2 d / p)$.

Definimos la \textbf{función de overhead total} $T_o$ como el trabajo que no realiza el algoritmo
secuencial:
\[
T_o(p) = p\,T_p - W = p\,T_{\text{comp}} + p\,T_{\text{comm}} - W.
\]
Como $p\,T_{\text{comp}} = W$, todo el overhead proviene de la comunicación y la sincronización:
\[
\boxed{\; T_o(p) = p\, T_{\text{comm}}. \;}
\]
Sustituyendo en la eficiencia se obtiene la forma canónica del curso:
\[
E = \frac{S_p}{p} = \frac{T_1}{p\,T_p} = \frac{W}{W + T_o}
  = \frac{1}{\,1 + \dfrac{T_o}{W}\,}.
\]
Para que $E=O(1)$ sea constante, el término de overhead $T_o/W$ debe permanecer constante: esa
igualdad, despejada, define la condición de isoeficiencia.

\subsubsection{Isoeficiencia en MPI (memoria distribuida)}

En MPI la comunicación es explícita y es el cuello de botella. Sustituyendo $n_{tr}=n_{te}=n$ en
$T_{\text{comm}}$ y conservando los términos dominantes:
\[
T_o = p\,T_{\text{comm}}
    \approx \underbrace{p^2\alpha}_{(1)\ \text{latencia}}
    \;+\; \underbrace{p\,\beta\, n\, d}_{(2)\ \text{scatter}}
    \;+\; \underbrace{p\,\beta\, n\, d\,\log p}_{(3)\ \text{bcast/reducción}}.
\]
El término (3) proviene de difundir el conjunto de prueba ($O(nd)$ bytes) y de reducir las listas
Top-$K$ a lo largo de una jerarquía de profundidad $\log p$; domina sobre (2). Imponemos
$T_o/W=\text{const}$ con $W=n^2 d$:
\[
(1)\quad \frac{p^2\alpha}{n^2 d}=\text{const} \;\Rightarrow\; n \propto p,
\qquad
(2)\quad \frac{p\,\beta n d}{n^2 d}=\frac{p\beta}{n}=\text{const} \;\Rightarrow\; n \propto p,
\]
\[
(3)\quad \frac{p\,\beta n d\,\log p}{n^2 d}=\frac{p\beta\log p}{n}=\text{const}
\;\Rightarrow\; n \propto p\,\log p.
\]
La condición más restrictiva es la del término (3). Por tanto:
\[
\boxed{\;
\text{Isoeficiencia MPI:}\quad n \propto p\,\log p,
\qquad
W_{iso} = \Theta\!\big(p^2 \log^2 p\big).
\;}
\]
Es decir, el dataset debe crecer ligeramente más rápido que $p$ (factor $\log p$) para amortizar
la difusión del test y la reducción global de candidatos sobre la red. Si la reducción se hiciera
con un \texttt{gather} plano al proceso raíz en vez de la reducción en árbol, el término dominante
sería $O(p)$ en lugar de $O(\log p)$, saturando el enlace del root y empeorando notablemente la
isoeficiencia.

\subsubsection{Isoeficiencia en OMP / PRAM (memoria compartida)}

En OpenMP los hilos comparten el mismo espacio de direcciones (modelo UMA), lo que cambia el
overhead de forma fundamental:
\begin{itemize}
    \item \textbf{No hay difusión de datos:} el conjunto de prueba es visible por todos los hilos,
    por lo que \emph{desaparecen} los términos $p\alpha$ y $\beta n d$ de la comunicación
    explícita ($T_{\text{comm}} \approx 0$).
    \item El único overhead proviene de la \textbf{sincronización de hilos} (barreras) y de la
    reducción de listas Top-$K$ de profundidad $O(\log p)$, realizada mediante accesos directos a
    memoria.
\end{itemize}
Modelando ese overhead (con $K$ constante):
\[
T_o^{\text{omp}} \approx p\cdot O(n_{te}K\log p) + O(p)
= O(p\,n\,\log p) + O(p).
\]
Imponiendo $T_o^{\text{omp}}/W=\text{const}$ con $W=n^2 d$, y observando que ya no aparece el
factor de transporte de datos $\beta d$:
\[
\frac{p\,n\,\log p}{n^2 d}=\frac{p\log p}{n\,d}=\text{const}
\;\Rightarrow\;
n \propto \frac{p\,\log p}{d}.
\]
\[
\boxed{\;
\text{Isoeficiencia OMP:}\quad n \propto \frac{p\,\log p}{d}\ \ (\text{a lo sumo } n\propto p),
\qquad
W_{iso} = \Theta(p^2)\ \text{sin la constante de red } \beta.
\;}
\]
OMP tiene una isoeficiencia asintóticamente mejor (no paga latencia ni ancho de banda de red),
\textbf{pero} $p$ está estrictamente acotado al número de núcleos físicos del nodo y todo el
dataset debe caber en su RAM.

\subsubsection{Comparación y elección del paradigma}

\begin{table}[H]
\centering
\begin{tabular}{lll}
\toprule
Criterio & MPI (memoria distribuida) & OMP (memoria compartida) \\
\midrule
Comunicación $T_{\text{comm}}$ & Explícita: $\alpha,\beta$ (red) & $\approx 0$ (acceso directo) \\
Overhead dominante & Difusión/reducción $\beta n d\log p$ & Sincronización (barreras) \\
Condición de isoeficiencia & $n \propto p\log p$ & $n \propto p\log p / d$ (más suave) \\
Trabajo $W_{iso}$ & $\Theta(p^2\log^2 p)$ & $\Theta(p^2)$ sin $\beta$ \\
Límite de $p$ & Horizontal (varios nodos) & Núcleos de un solo nodo \\
Límite de $n$ & Suma de RAM de los nodos & RAM de un solo nodo \\
\bottomrule
\end{tabular}
\caption{Comparación de la escalabilidad de ambos paradigmas para el KNN.}
\label{tab:iso-comparacion}
\end{table}

Aunque OMP escala mejor asintóticamente, su paralelismo está limitado a los núcleos y la memoria
de un solo nodo. En este proyecto se eligió \textbf{MPI por un tema de memoria}: al distribuir el
conjunto de entrenamiento, el mismo código puede superar el límite de RAM de un nodo y escalar
$n$ y $p$ horizontalmente sin reescritura. Los experimentos evaluaron hasta $p=32$ procesos, y la
elección del paradigma de memoria distribuida es la que habilita crecer más allá de ese punto
cuando el dataset no cabe en un único nodo.



\section{Conclusiones}

El diseño paralelo propuesto aprovecha que el cálculo de distancias en KNN es altamente
paralelizable. Al distribuir el conjunto de entrenamiento, el costo dominante se reduce de:
\[
O(n_{tr}n_{te}(d+\log K))
\]
a:
\[
O\left(
\frac{n_{tr}n_{te}}{p}(d+\log K)
\right),
\]
más costos adicionales de comunicación y reducción.

La reducción binaria de Top-$K$ evita enviar todos los resultados completos al proceso raíz
mediante una comunicación plana, reduciendo la profundidad de combinación a:
\[
O(\log p).
\]

Se espera buen speedup para valores moderados de $p$, especialmente cuando:
\[
n_{tr}n_{te}d
\]
es suficientemente grande para amortiguar los costos de comunicación. Sin embargo, para valores
altos de $p$, la eficiencia disminuye debido a la difusión del conjunto de prueba, la distribución
del entrenamiento y la reducción global de candidatos.

\begin{thebibliography}{9}

\bibitem{mpi4py}
MPI for Python documentation.
\textit{mpi4py}. Disponible en la documentación oficial de mpi4py.

\bibitem{sklearn}
Scikit-learn documentation.
\textit{The digits dataset}. Disponible en la documentación oficial de scikit-learn.

\bibitem{course}
José Fiestas.
\textit{Computación Paralela y Distribuida: Proyecto de Laboratorio 2026-I}.
Universidad de Ingeniería y Tecnología, 2026.

\end{thebibliography}

\end{document}